\documentclass{article}
\usepackage[utf8]{inputenc}
\usepackage[spanish]{babel}

\title{Memoria Nexus Gaming Club}
\author{Albert y Sergio}
\date{Curso 2026-2027}

\begin{document}

\maketitle

\section{Introducción, problema y objetivos}

\subsection{Nombre del proyecto y frase de venta}
Nuestro proyecto se llama Nexus Gaming Club. 
La frase de venta es: Infraestructura de red y sistemas de alto rendimiento para un edificio físico dedicado al gaming y eSports.

\subsection{Problema o necesidad}
Los clubs de videojuegos tienen varios problemas hoy en día:
\begin{itemize}
    \item No tienen red propia y dependen de plataformas externas.
    \item La red no está preparada para jugar con buen ping.
    \item No hay servidores propios de juego ni de voz.
    \item La información está repartida por varios sitios.
    \item No hay seguridad real (ni firewall, ni copias de seguridad).
\end{itemize}

\subsection{Usuarios}
Los usuarios que van a usar nuestro sistema son:
\begin{itemize}
    \item Administradores: Para gestionar usuarios y servidores.
    \item Moderadores: Para gestionar el contenido.
    \item Jugadores: Para acceder a los escritorios y torneos.
        \item Visitantes: Para ver la página web.
\end{itemize}

\subsection{Objetivos}
Los objetivos que queremos conseguir este curso son:
\begin{enumerate}
    \item Montar la red local y el wifi.
    \item Configurar los servidores (DHCP, DNS, etc.).
    \item Hacer una página web con HTML y CSS.
    \item Poner un firewall y hacer copias de seguridad.
    \item Documentar todo en latex y subirlo a GitHub.
\end{enumerate}

\end{document}